\cleardoublepage

\chapter{Introducción}
\label{cap:intro} % etiqueta para poder referenciar luego en el texto con ~\ref{sec:intro}
\pagenumbering{arabic} % para empezar la numeración de página con números

\section{Contexto}
\label{sec:cont}



\section{Estructura de la Memoria}
\label{sec:estr}

A continuación, se describe la estructura de esta memoria. Está organizada en 7 capítulos que abarcan el desarrollo completo de este Trabajo Fin de Grado:

\begin{itemize}
    \item \textbf{Capítulo 1 - Introducción:} Se presenta de forma concisa el pensamiento computacional y su relación con \textit{Snap!} y \textit{Dr. Snap!} en el ámbito educativo. Asimismo, se expone la motivación del proyecto y se describe brevemente su estructura.
    \item \textbf{Capítulo 2 - Objetivos:} Se detallan los objetivos generales y específicos que se pretenden alcanzar con el desarrollo de este proyecto; así como se presenta una planificación temporal de las tareas a realizar para alcanzarlos.
    \item \textbf{Capítulo 3 - Estado del Arte y Tecnologías Utilizadas:} Se realiza un estudio sobre las bases del pensamiento computacional; además de analizar el funcionamiento de \textit{Snap!} y la existencia de herramientas para el análisis de proyectos realizados con este lenguaje. Asimismo, se estudian los predecesores de \textit{Dr. Snap!} y sus versiones anteriores. También se describen los lenguajes de programación, marcado y \textit{frameworks} empleados para el desarrollo de este proyecto.
    \item \textbf{Capítulo 4 - Diseño e Implementación:} Se describen las mejoras realizadas en \textit{Dr. Snap!}; así como el desarrollo del Modo de Análisis por Lotes (\textit{Batch Mode}), de nuevos \textit{Malos Hábitos} a analizar como Código Duplicado y el despliegue de una \textit{versión beta} en un servidor universitario.
    \item \textbf{Capítulo 5 - Uso en el Ámbito de la Investigación:} Se exponen los resultados de un Análisis Empírico y Evaluación de Límites Lógicos en \textit{Snap!} realizado con la \textit{versión beta} desarrollada.
    \item \textbf{Capítulo 6 - Casos de Uso del Sistema:} Se describen tres casos de uso\footnote{Hay dos casos de uso en los que el actor principal es el profesor y otro más en el que lo es el alumno.} de la aplicación web desarrollada, exponiendo su funcionamiento y justificando su ayuda en el desarrollo del pensamiento computacional con ejemplos reales de dichos casos de uso.
    \item \textbf{Capítulo 7 - Conclusiones:} Se exponen los resultados del Trabajo valorando la consecución de los objetivos, aplicación de lo aprendido y lecciones aprendidas. Además, se presentan futuras vías de desarrollo de la aplicación web.
\end{itemize}
